\section{Descrição do grafo original}
\label{sec:descricao}

O web-flickr faz parte da coleção de grafos do SNAP (Stanford Large Network Dataset Collection)~\cite{snapnets} e foi construído por McAuley e Leskovec~\cite{mcauley2012image} para estudar a rotulação de imagens com o uso de metadados de redes sociais.

Cada vértice é uma imagem do Flickr. Uma aresta liga duas imagens que compartilham metadados: imagens do mesmo local, submetidas à mesma galeria, grupo ou conjunto, com tags em comum, do mesmo usuário ou tiradas por amigos. A relação é simétrica, então o grafo é não direcionado. O SNAP distribui o grafo como uma lista de arestas comprimida, \texttt{flickrEdges.txt.gz}, em que cada par não ordenado de vértices aparece uma vez. Além dela, o SNAP distribui atributos dos vértices, atributos das arestas e os metadados das imagens em XML. Este trabalho usa só a lista de arestas, e a Seção~\ref{sec:discussao} explica por que os atributos das arestas não foram usados como pesos.

O SNAP declara as seguintes estatísticas para o grafo~\cite{snapnets}: 105.938 vértices, 2.316.948 arestas, maior componente conexa com 105.722 vértices, coeficiente de clusterização médio de 0,0891, 107.987.357 triângulos e diâmetro 9. Esses valores servem apenas de referência: o tratamento dos dados é descrito na Seção~\ref{sec:tratamento}, e as medidas são recalculadas a partir do arquivo original na Seção~\ref{sec:parte1}.

\subsection{Tratamento realizado}
\label{sec:tratamento}

O arquivo original foi convertido em um grafo simples não direcionado, e dele se extraiu a maior componente conexa. A Tabela~\ref{tab:tratamento} compara os três conjuntos.

\begin{table}[htbp]
  \centering
  \caption{Vértices, arestas, auto-laços e arestas repetidas da lista de arestas original, do grafo simples obtido pelo tratamento e da sua maior componente conexa.}
  \label{tab:tratamento}
  \begin{tabular}{lrrr}
\toprule
 & Lista bruta & Grafo simples & Maior componente \\
\midrule
Vértices & 105.938 & 105.938 & 105.722 \\
Arestas & 2.316.948 & 2.316.948 & 2.316.668 \\
Auto-laços & 0 & 0 & 0 \\
Arestas repetidas & 0 & 0 & 0 \\
\bottomrule
\end{tabular}

\end{table}

Cada linha do arquivo foi tratada como uma aresta não ordenada, porque a relação entre imagens que compartilham metadados é simétrica. Os auto-laços foram removidos e, entre as arestas que ligam o mesmo par de vértices, em qualquer ordem, manteve-se uma só, de modo que o resultado é um grafo simples. A Tabela~\ref{tab:tratamento} mostra que o arquivo original não tinha auto-laços nem arestas repetidas, o que confirma a declaração do SNAP de que cada par aparece uma vez: o grafo simples tem os mesmos vértices e arestas da lista original. O conjunto de vértices é o dos extremos das arestas, pois o arquivo não lista vértices isolados.

Os vértices foram renumerados de $0$ a $n-1$, na ordem crescente do identificador original, para que possam indexar diretamente as listas de adjacência usadas pelas implementações. O identificador original de cada vértice foi guardado em um mapa, que permite voltar às imagens do Flickr.

O grafo é desconexo: a maior componente conexa não contém todos os vértices (Tabela~\ref{tab:tratamento}). O seu tamanho coincide com os 105.722 vértices declarados pelo SNAP. A maior componente recebeu numeração própria, de $0$ a $n_{\mathrm{lcc}}-1$, também registrada no mapa de identificadores.

A política adotada para o grafo desconexo foi a seguinte:
\begin{itemize}
  \item O grafo simples completo é usado nas medidas que não dependem de conexidade (número de vértices e de arestas, graus e sua distribuição, densidade, componentes conexas, coeficiente de clusterização e triângulos) e nos algoritmos que se aplicam a grafos desconexos (busca em largura e em profundidade, verificação de eulerianidade, pontos de articulação e pontes, Dijkstra).
  \item A maior componente conexa é usada nas medidas de distância (diâmetro, raio, comprimento médio dos caminhos e o comprimento médio da análise small-world), porque a distância entre vértices de componentes diferentes é infinita. Ela também é usada na árvore geradora mínima, porque só um grafo conexo tem árvore geradora, e na análise de robustez, como recomenda a descrição do trabalho.
\end{itemize}
Como a maior componente reúne quase todos os vértices e arestas do grafo (Tabela~\ref{tab:tratamento}), as análises feitas sobre ela representam o grafo quase inteiro.
